\documentclass[a4paper,fontsize=11pt,fleqn]{scrartcl}%fleqn pour aligner les equations à gauche

\usepackage{../../Styles/Eval}

\newcommand{\chnum}{Chapitre 5 et 8}
\newcommand{\chtitle}{Polarité des molécules et titrages colorimétriques}
\newcommand{\dtype}{DS5}

\begin{document}
\maketitle

\section{Polarité des molécules organiques}

On considère les structures chimiques suivantes : \ce{HSO4-} ; \ce{HPO3} et \ce{H2CO2}
\begin{enumerate}
    \item Proposer pour chaque molécule un schéma de Lewis correct.
    \item En se basant sur le tableau des électronégativités de Pauling, réécrire les structures de Lewis en indiquant les éventuelles charges partielles qui y sont présentes.
    \item Déterminer la géométrie de ces structures autour des atomes de soufre, de phosphore et du carbone.
    \item Déterminer si chacune des molécules est polaires ou apolaires.
\end{enumerate}

\section{Dosage de l'eau de Javel}
L'eau de Javel est une solution oxydante utilisée comme désinfectant et comme décolorant. Selon l'emploi désiré, il est important de connaître sa concentration.

\begin{documentboxenv}{Synthèse de l'eau de Javel}
Charles Guillaume Scheele (1742-1786), pharmacien suédois, découvrit le dichlore au 18\textsuperscript{ème} siècle. Ce gaz intervient dans la fabrication de l'eau de Javel ;  qui doit son nom à un ancien village aujourd'hui compris dans un quartier de Paris.\\
C'est à Javel que Claude Louis Berthollet (1748-1822), directeur à la manufacture des Gobelins, fabriqua ce produit décolorant et désinfectant et l'employa en 1785 pour le blanchiment des toiles textiles. L'eau de Javel est obtenue par action du dichlore \ce{Cl2(g)} sur l'hydroxyde de sodium ou soude et contient des ions \ce{Cl-(aq)}., \ce{ClO-(aq)}, \ce{Na+(aq)} ainsi que de l'eau \ce{H2O l)}.\\
C'est donc une solution aqueuse constituée entre autres d'ions chlorure et d'ions hypochlorite \ce{ClO-(aq)}. En milieu acide, l'eau de Javel subit une transformation complète représentée par la réaction d'équation : $$\ce{ClO-(aq) + Cl-(aq) + 2 H+(aq) -> H2O(l) + Cl2(g)} $$
Cette transformation permet de définir le degré chlorométrique. Celui-ci est égal au volume, exprimé en litres, de dichlore produit par un litre d'eau de Javel. Ce volume est mesuré à une température de \SI{0}{\celsius} sous une pression de \SI{1,013}{\bar}.
\end{documentboxenv}
\begin{documentboxenv}{Méthode expérimentale}
Pour vérifier l'indication portée sur une bouteille commerciale d'eau de Javel, \SI{12}{\degree chl} (12 degrés chlorométriques), on procède à une expérience.\\
\textbf{Principe de la manipulation}
\begin{itemize}
    \item On ajoute un excès d'ions iodure à un volume connu de solution d'eau de Javel
    \item En milieu acide (concentré en \ce{H+}), les ions hypochlorite \ce{ClO-(aq)} oxydent les ions iodure \ce{I-(aq)}
    \item L'équation de la réaction est : $$ \ce{ClO-(aq) + 2 I-(aq) + 2H+(aq) -> H2O(l)  + I2(g) + Cl-(aq)} \hspace{2em} (1)$$
\end{itemize}
    On considère cette réaction comme totale. Le diiode formé appartenant au couple \ce{I2(aq)}/\ce{I-(aq)} est quantifié en le faisant réagir avec les ions thiosulfate, réducteurs du couple \ce{S4O6^2-(aq)}/\ce{S2O3^2-(aq)}. On en déduit alors la quantité d'ions hypochlorite puis le degré chlorométrique.
\end{documentboxenv}

    \subsection{Réaction entre les ions hypochlorite et les ions iodure}
    \begin{enumerate}
        \item L'eau de Javel commerciale étant trop concentrée, il faut  d'abord effectuer une dilution du dixième pour obtenir \SI{50,0}{mL} de solution diluée $S$. Décrire une méthode qui permet d'effectuer cette dilution. Préciser la verrerie nécessaire (noms et volumes).
        \item Dans un erlenmeyer (noté $E$), on introduit dans cet ordre : $V = \SI{10,0}{mL}$ de solution $S$ puis $V' = \SI{20,0}{mL}$ de la solution d'iodure de potassium (\ce{K+(aq)} ; \ce{I-(aq)}). Indiquer la verrerie à utiliser.
    \end{enumerate}

    \subsection{Détermination de la concentration en eau de Javel}
    
        \begin{documentboxenv}{Montage pour dosage}
        \noindent\begin{minipage}{.49\linewidth}
        \begin{center}
            \includegraphics[width = .8\linewidth,trim={0 2cm 20cm 1.5cm},clip]{Ressources/Capture d’écran 2026-03-28 à 20.17.49.png}
        \end{center}
    \end{minipage}
    \begin{minipage}{.49\linewidth}
        \textbf{Procédure :} À l'aide d'une burette graduée, on verse peu à peu, dans l'erlenmeyer $E$, une solution de thiosulfate de sodium de formule (\ce{2Na+(aq)} ; \ce{S2O3^2-(aq)}) de concentration molaire apportée $c_1 = \SI{0,10}{\mole\per\liter}$. L'ajout se fait sous agitation à l'aide d'un agitateur magnétique. On ajoute une pointe de spatule d'empois d'amidon afin de mieux repérer le changement de couleur. Le volume de thiosulfate pour lequel le mélange dans l'erlenmeyer est décoloré totalement est $V_{1E} = \SI{10,0}{mL}$.\\
        \textit{En présence de diiode en solution, l'empois d'amidon prend une couleur bleue foncée.}
    \end{minipage}
    \end{documentboxenv}
    
    
    \begin{enumerate}
        \item Écrire l'équation de la réaction qui sera notée (2), entre le diiode et les ions thiosulfate.
        \item Justifier la décoloration progressive du mélange réactionnel.
        \item Quel est le réactif limitant avant l'équivalence ? Après l'équivalence ? À l'équivalence ?
        \item Déduire des résultats de l'expérience, la quantité de matière de diiode présente dans le mélange réactionnel. Cette quantité de matière correspond aussi à la quantité produite lors de la réaction (1).
        \item Calculer la quantité de matière d'ions hypochlorite initialement présents dans le prélèvement de volume $V$.
        \item Déterminer la concentration en ions hypochlorite de la solution $S$, puis de la solution commerciale.
        \item En utilisant l'équation de la réaction chimique donnée au Doc. 1, calculer la quantité de matière de dichlore produite par un litre d'eau de Javel.
        \item Le volume molaire d'un gaz parfait, dans les conditions de température et de pression citées dans le texte vaut \SI{22,4}{\liter\per\mole}. En déduire le degré chlorométrique de l'eau de Javel commerciale utilisée.
    \end{enumerate}
\end{document}

